\chapter{Boundary connector: минимальная структурная спецификация}

Новый boundary action должен решить три задачи из одного источника:
разрушить family degeneracy, выбрать CP orientation и создать portal между
Version III.H и observed sector. Чтобы это не превратилось в три независимых
коэффициента, сначала классифицируем минимальные connector packages по
структурным возможностям.

\section{Обязательные свойства}

Кандидат проходит первый boundary gate только если одновременно:
\begin{enumerate}
\item содержит off-diagonal finite edge между hidden и observed blocks;
\item допускает nonzero portal, возникающий из того же connector orbit;
\item несёт family-triplet datum, способный действовать на
\(\operatorname{im}P_3\);
\item содержит closed oriented path с rephasing-invariant phase;
\item является anomaly-free до интегрирования тяжёлых messengers.
\end{enumerate}

\section{Restricted connector menu}

\begin{center}
\small
\begin{tabular}{>{\raggedright\arraybackslash}p{0.25\textwidth}ccccc}
\toprule
Package & Edge & Portal & Family & CP loop & Anomaly \\
\midrule
Direct sum & no & no & no & no & yes \\
Один scalar triplet & no & yes & yes & no & yes \\
Одна vectorlike pair & yes & yes & no & no & yes \\
Triplet + одна messenger chain & yes & yes & yes & no & yes \\
Two-path vectorlike triplet cycle & yes & yes & yes & yes & yes \\
\bottomrule
\end{tabular}
\end{center}

\begin{theorem}[Single-connector no-go в restricted menu]
Ни один single-field orbit и ни одна открытая messenger chain не могут
одновременно породить portal, family breaking и физическую CP phase. Scalar
triplet не создаёт off-diagonal spectral bridge; одна vectorlike pair не
несёт family carrier; одна chain не содержит gauge-invariant closed phase и
остаётся rephasing-real.
\end{theorem}

Первый package, проходящий capability gate, содержит две vectorlike messenger
paths и один family-triplet connector. Схематически:
\begin{equation}
\Psi_{\rm obs}
\longrightarrow X_1
\mathrel{\substack{\longrightarrow\\[-0.6ex]\longrightarrow}}
X_2
\longrightarrow\Psi_{\rm obs},
\end{equation}
где две средние arrows являются неэквивалентными family/orientation edges, а
хотя бы один промежуточный orbit несёт hidden charge. Замкнутая разность двух
paths может содержать rephasing-invariant Wilson/Krajewski phase.

\section{Статус минимальности}

Полученный минимум относится только к объявленному renormalizable connector
menu. Он не доказывает существование допустимой real finite geometry.
Следующие условия ещё обязательны:
\begin{enumerate}
\item совместимость с \(J\), grading и order-one condition;
\item явная таблица SM representations и hidden charges;
\item cancellation всех continuous и discrete anomalies;
\item доказательство, что loop phase не удаляется field redefinitions;
\item вывод relative coefficients portal/family/CP из одного spectral trace;
\item отсутствие произвольной \(M_3\)-матрицы после интегрирования messengers.
\end{enumerate}

\begin{nogocriterion}[Запрет capability-to-existence jump]
Boolean capability pass не является построенным boundary action. Пока не
предъявлены representation, finite Dirac operator и anomaly ledger, two-path
cycle считается только первым admissible graph ansatz.
\end{nogocriterion}

Машинный capability ledger сохранён в
\path{s2t_v4_boundary_connector_menu_gate_results.json}.